\documentclass[a4paper,11pt]{article}
\usepackage{hyperref}

\begin{document}
\title{Preface by the PHOCON 2026 Chairs}
\date{}
\maketitle

The \emph{International Workshop on
  Petri Nets, Higher-Dimensional Automata, Partial Orders, and Concurrency}
(\emph{PHOCON})
took place on 23 June 2026.
It was co-located with the 47th International Conference on Application and Theory of Petri Nets
(PETRI NETS 2026)
in Hamburg, Germany.

PHOCON is a workshop on concurrency theory, in particular Petri nets and non-interleaving concurrency. It concerns itself with models for concurrency, partial-order semantics, their relation to Petri nets, and their applications. Topics of interest include, but are not limited to
\begin{itemize}
\item event structures,
\item configuration structures,
\item negotiations,
\item higher-dimensional automata and variants,
\item partial-order semantics,
\item unfoldings,
\item partial-order reduction,
\item category theory,
\item history-preserving bisimulation and variants,
\item real-time concurrency.
\end{itemize}
Further information is available at \url{https://phocon.github.io/}.

PHOCON 2026 received one regular and five abstract-only submissions.
Each submission was reviewed by three referees, and based on the reviews, all submissions were accepted.
Additionally, the program of PHOCON 2026 included invited talks by \emph{Stefan Haar} and \emph{Rob van Glabbeek}.
We wish to thank the members of the program committee, the invited speakers,
and the local organisers in Hamburg for their support.

\bigskip\noindent
June 2026
\hfill
\begin{minipage}[t]{.5\linewidth}
  \flushright
  Uli Fahrenberg \\
  Loïc Hélouët \\
  Philipp Schlehuber-Caissier \\
  Krzysztof Ziemiański
\end{minipage}

\end{document}
